% linear_chain2.py — описание для отчёта научному руководителю.
% Сборка: pdflatex linear_chain2.tex (рекомендуется texlive-lang-cyrillic).

\documentclass[a4paper,11pt]{article}

\usepackage[T2A]{fontenc}
\usepackage[utf8]{inputenc}
\usepackage[russian]{babel}
\usepackage{amsmath,amssymb}
\usepackage{geometry}
\geometry{left=2.5cm,right=2.5cm,top=2.5cm,bottom=2.5cm}
\usepackage{hyperref}
\usepackage{enumitem}
\usepackage{array}
\setlist{nosep,leftmargin=*}

\hypersetup{colorlinks=true,linkcolor=blue,urlcolor=blue}

\newcommand{\code}[1]{\texttt{\detokenize{#1}}}

\title{Сценарий \code{linear_chain2.py} модуля\\ \texttt{Drone\_Swarm\_Simulator\_v2}:\\
математическая модель, реализация и обоснование решений}
\author{}
\date{}

\begin{document}
\maketitle

\begin{abstract}
В отчёте рассматривается Python-сценарий \code{scenarios/linear_chain2.py}, предназначенный для численного эксперимента в связке нескольких экземпляров ArduPilot SITL. Цель --- сформировать \textbf{линейную цепочку} беспилотных аппаратов в горизонтальной плоскости: крайние агенты (якоря) удерживают концы отрезка~$AB$, внутренние стабилизируют порядок на линии. Управление реализовано через переопределение радиоканалов (roll/pitch) в режиме позиционного удержания. Внутренний закон \textbf{гибридный}: геометрическая компонента (середина между соседями по проекции на~$AB$ и удержание на прямой) суммируется с уменьшенным вкладом распределённого закона из постановки «коридор'' (\texttt{corridor\_sweeping}): учёт ближайших соседей в четырёх полупространствах, нелинейность $\tanh$, виртуальный поперечный зазор при отсутствии наблюдения.

Ниже последовательно изложены постановка задачи, программная архитектура, системы координат, математические соотношения с пояснением каждой вновь вводимой величины, соответствие функций исходному коду и мотивы инженерных компромиссов.
\end{abstract}

\tableofcontents
\newpage

%====================================================================
\section{Введение и постановка задачи}

\subsection{Цель работы}
Требуется организовать совместное поведение роя из $N$ беспилотных летательных аппаратов (в коде $N\ge 4$) в горизонтальной плоскости так, чтобы в установившемся режиме:
\begin{itemize}
  \item два агента с \textbf{минимальным} и \textbf{максимальным} целочисленным идентификатором \code{id} выступали \emph{якорями} и занимали (с допустимой погрешностью) заранее заданные концы отрезка~$AB$ длины~$L$;
  \item остальные агенты располагались на той же прямой между якорями с \textbf{близким к равномерным} упорядочением по положению вдоль отрезка.
\end{itemize}

Здесь $N$ --- число дронов в эксперименте; $L$ --- заданная длина отрезка в метрах (параметр \code{--segment-length}, по умолчанию совпадает с константой \code{SEGMENT\_LENGTH\_M} в файле).

\subsection{Ограничения среды моделирования}
Эксперимент выполняется в программной связке с ArduPilot SITL. На уровне сценария нет прямого задания вектора тяги в инерциальном базисе; вместо этого используется режим, эквивалентный удержанию позиции (\code{POS\_HOLD}), и низкоуровневое \textbf{переопределение каналов RC} (\code{MAV\_OVERRIDE}). Следовательно, закон управления должен быть сформулирован как \emph{желаемое горизонтальное смещение} (в плоскости North--East), затем переведён в оси корпуса (вперёд--вправо) и дискретизован в приращения ширины импульса PWM по каналам крена и тангажа.

\subsection{Отличие от сценария \code{linear_chain.py}}
В версии~1 (\code{linear_chain.py}) внутренний дрон строит двумерную целевую точку на отрезке (середина между соседями по одномерной координате вдоль~$AB$), сглаживает её и подаёт на PID-контур по roll/pitch. В \code{linear_chain2.py} \textbf{явная 2D-цель для внутренних не вводится}: формируется вектор $\mathbf{c}_{xy}$ в плоскости North--East как сумма геометрической и «коридорной'' компонент, после чего к полученным сигналам в корпусе применяется пропорциональное преобразование в PWM с ограничениями, минимальным шагом и ограничением скорости изменения команды (slew rate). Такой подход проще согласовать с логикой \texttt{corridor\_sweeping} и позволяет отдельно настраивать «жёсткость'' выравнивания по линии и «мягкость'' нелинейного распределённого вклада.

%====================================================================
\section{Архитектура программного модуля}

\subsection{Общая схема}
Модуль запускает несколько потоков POSIX: для каждого дрона --- инициализация и основной контур поведения; дополнительно --- поток обмена координатами между логическими контроллерами. Синхронизация старта после взлёта обеспечивается объектом \code{threading.Barrier}.

\subsection{Инициализация агента}
Функция \code{initialize\_drone\_parallel} для каждого экземпляра \code{DroneController} выполняет подключение по MAVLink, последовательность \code{INIT\_STEPS}: перевод в безопасный режим, армирование, взлёт на высоту \code{TAKEOFF\_ALT\_M}, переход в \code{POS\_HOLD}, установка нейтральных значений RC, запрос стримов позиции и ориентации с заданной частотой, запуск периодической посылки keepalive для удержания override. Таким образом, к началу эксперимента все агенты находятся в сопоставимом автопилотном режиме.

\subsection{Обмен информацией о положении}
Поток \code{coordinate\_exchange\_loop} с частотой $f_{\mathrm{ex}}$ (аргумент \code{--exchange-hz}, величина $f_{\mathrm{ex}}$ измеряется в герцах) опрашивает локальные оценки положения, переводит их в \emph{общий кадр} (см.\ раздел~\ref{sec:coords}) и записывает позиции остальных агентов в память каждого контроллера (\code{update\_other\_drone\_position}). Важно: это \textbf{не} модель бортовой связи; глобальная сводка позиций формируется в одном процессе симулятора. Для распределённого алгоритма на реальном рое потребовалась бы замена этого блока протоколом обмена измерениями.

\subsection{Контуры управления по ролям}
После фиксации концов отрезка $A$ и $B$ (см.\ раздел~\ref{sec:segment}) для якорей запускается \code{endpoint\_loop}, для внутренних --- \code{internal\_chain\_anatoliy\_loop}. Период дискретизации контура управления $T_c = 1/f_c$, где $f_c$ --- \code{CONTROL\_HZ} (в текущей версии файла $f_c=18$\,Гц). Именно с шагом $T_c$ вычисляются команды RC для соответствующего агента.

\subsection{Точка входа}
Функция \code{main()} разбирает аргументы командной строки, создаёт список конфигураций (идентификатор, UDP-порт и т.\,д.), поднимает потоки и по завершении корректно останавливает контроллеры. Параметры, влияющие на математику внутреннего закона: \code{--r-vis} (радиус видимости соседей $R_{\mathrm{vis}}$), \code{--w-cross} (виртуальный поперечный зазор $w_{\mathrm{cross}}$).

%====================================================================
\section{Системы координат и отрезок цепи}
\label{sec:coords}

\subsection{Горизонтальная плоскость NED и поля позиции в коде}
В авиационной системе North--East--Down (NED) горизонтальные оси: \textbf{North} (на север) и \textbf{East} (на восток). В словаре позиции сценария компоненты обозначены \code{x} и \code{y}; в дальнейшем отождествляем
\begin{center}
\code{x} $\equiv$ смещение на север ($N$),\quad \code{y} $\equiv$ смещение на восток ($E$).
\end{center}
Компонента высоты \code{z} в NED направлена вниз; целевая высота взлёта задаётся отрицательным значением \code{-TAKEOFF\_ALT\_M} в тех же единицах, что ожидает автопилот.

\subsection{Общий кадр (common frame)}
У разных экземпляров SITL локальные координаты смещены относительно друг друга. Чтобы сравнивать положения агентов в одной системе, вводится поправка только по оси East (поле~\code{y}):
\begin{equation}
\label{eq:ycommon}
y_{\mathrm{common}} = y_{\mathrm{local}} + \Delta y(\mathrm{id}).
\end{equation}
Здесь $y_{\mathrm{local}}$ --- восточная координата, возвращаемая API контроллера для данного дрона; $\mathrm{id}\in\{1,2,\ldots\}$ --- целочисленный идентификатор агента в конфигурации сценария; $\Delta y(\mathrm{id})$ --- детерминированный сдвиг. В реализации $\Delta y(\mathrm{id}) = (\mathrm{id}-1)\cdot 2$\,м (функция \code{\_did\_offset\_y}). Смысл: развести траектории на визуализации и уменьшить ложное «наложение'' траекторий при чтении локальных координат разных SITL. Северная координата \code{x} в common frame не сдвигается.

\subsection{Построение концов отрезка $A$ и $B$}
\label{sec:segment}
Пусть после взлёта в common frame измерены положения левого и правого якоря; их восточные координаты обозначим $y_{\mathrm{left}}$ и $y_{\mathrm{right}}$ (это значения \code{y} после преобразования~\eqref{eq:ycommon} для якорей с минимальным и максимальным \code{id}). Задаётся средняя ордината
\begin{equation}
\label{eq:ymid}
y_{\mathrm{mid}} = \frac{y_{\mathrm{left}} + y_{\mathrm{right}}}{2}.
\end{equation}
Величина $y_{\mathrm{mid}}$ --- скаляр, фиксирующий \textbf{одну} горизонтальную прямую для геометрии цепи: оба конца отрезка лежат на одной «высоте'' по East, даже если сами якоря в момент снимка слегка разошлись по~$y$ из-за шума оценки или задержек взлёта.

Длина отрезка $L>0$ задаётся параметром запуска. Точки $A$ и $B$ в горизонтальной плоскости (в координатах North--East относительно принятого начала common frame) задаются как
\begin{equation}
\label{eq:AB}
A = \bigl(N_A,\, E_A\bigr) = \bigl(-\tfrac{L}{2},\, y_{\mathrm{mid}}\bigr),\qquad
B = \bigl(N_B,\, E_B\bigr) = \bigl(+\tfrac{L}{2},\, y_{\mathrm{mid}}\bigr).
\end{equation}
Таким образом, $N_A,N_B$ --- северные координаты концов; $E_A=E_B=y_{\mathrm{mid}}$ --- восточные. В коде тем же значениям соответствуют поля \code{x} и \code{y} словарей \code{endpoint\_a\_common} и \code{endpoint\_b\_common}.

\subsection{Ортонормированный базис вдоль цепи}
Введём единичный вектор направления от $A$ к $B$ в плоскости North--East:
\begin{equation}
\label{eq:u}
\mathbf{u} = \frac{(N_B-N_A,\; E_B-E_A)^\top}{\sqrt{(N_B-N_A)^2 + (E_B-E_A)^2}} = (u_N,\, u_E)^\top .
\end{equation}
Компоненты $u_N$ и $u_E$ --- направляющие косинусы оси цепи. Левый ортогональный вектор в той же плоскости (поворот на $+90^\circ$ против часовой стрелки в плане North--East, если смотреть сверху вниз в стандартной ориентации NED):
\begin{equation}
\label{eq:n}
\mathbf{n} = (-u_E,\, u_N)^\top = (n_N,\, n_E)^\top .
\end{equation}
Вектор $\mathbf{u}$ направлен вдоль цепи от левого якоря к правому; $\mathbf{n}$ задаёт ось «влево'' от направления полёта вдоль цепи. Реализация: \code{\_unit\_vector\_xy}, \code{\_normal\_left\_xy}.

\subsection{Одномерная координата вдоль отрезка}
Для произвольной точки $p$ с горизонтальными координатами $(N_p,E_p)$ определим скалярную координату вдоль отрезка (проекцию на ось от точки~$A$):
\begin{equation}
\label{eq:s}
s(p) = (N_p - N_A)\, u_N + (E_p - E_A)\, u_E .
\end{equation}
Величина $s(p)$ измеряется в метрах; она монотонно возрастает при движении от~$A$ к~$B$ вдоль прямой. Эта координата используется для упорядочивания агентов «вдоль цепи'' независимо от наклона прямой в плоскости (при невырожденном $A\neq B$).

%====================================================================
\section{Математическая модель внутреннего контура}
\label{sec:internal}

Внутренний агент (не якорь) в каждый момент дискретного времени с шагом $T_c$ выполняет вычисления ниже. Обозначим через $p_{\mathrm{me}}$ вектор его положения в common frame в горизонтальной плоскости $(N_{\mathrm{me}},E_{\mathrm{me}})$, через $\mathcal{P}$ --- множество идентификаторов агентов, чьи позиции доступны (своя и чужие после обмена).

\subsection{Множество наблюдаемых соседей и радиус видимости}
Пусть $R_{\mathrm{vis}}>0$ --- радиус видимости (параметр \code{R\_VIS\_M} или \code{--r-vis}). Для каждого другого агента $j\in\mathcal{P}$, $j\neq \mathrm{my\_id}$, вычисляется евклидово расстояние в плоскости
\begin{equation}
d_j = \bigl\| p_{j,xy} - p_{\mathrm{me},xy} \bigr\|_2 .
\end{equation}
Здесь $p_{j,xy}$ --- горизонтальные координаты агента~$j$; $\mathrm{my\_id}$ --- идентификатор рассматриваемого дрона. Агент $j$ учитывается, если $d_j \le R_{\mathrm{vis}}$ и $d_j$ не ниже малый порог (в коде исключение нулевого вектора). Для каждого учтённого $j$ вводится вектор относительного положения в плоскости
\begin{equation}
\label{eq:Delta}
\boldsymbol{\Delta}_j = p_{j,xy} - p_{\mathrm{me},xy} = (\Delta_{j,N},\, \Delta_{j,E})^\top .
\end{equation}
Проекции $\boldsymbol{\Delta}_j$ на базис цепи $(\mathbf{u},\mathbf{n})$:
\begin{equation}
\label{eq:ajcj}
a_j = \boldsymbol{\Delta}_j\cdot \mathbf{u} = \Delta_{j,N} u_N + \Delta_{j,E} u_E,\qquad
c_j = \boldsymbol{\Delta}_j\cdot \mathbf{n} = \Delta_{j,N} n_N + \Delta_{j,E} n_E .
\end{equation}
Скаляр $a_j$ --- смещение соседа \textbf{вдоль} цепи относительно рассматриваемого дрона (положительное --- в сторону $\mathbf{u}$); $c_j$ --- смещение \textbf{поперёк} (положительное --- в сторону $\mathbf{n}$). Если число таких соседей меньше двух, команда RC сбрасывается в нейтраль (недостаточно информации для устойчивого закона).

\subsection{Четыре полупространства: ближайшие расстояния}
По совокупности пар $(a_j,c_j)$ находятся минимальные положительные смещения в каждом из четырёх квадрантов осей $(a,c)$. Вводим:
\begin{align*}
d_{a+} &= \min\{ a_j : a_j\ge 0,\; j\ \text{учтён} \}, &
d_{a-} &= \min\{ |a_j| : a_j< 0,\; j\ \text{учтён} \},\\
d_{c+} &= \min\{ c_j : c_j\ge 0,\; j\ \text{учтён} \}, &
d_{c-} &= \min\{ |c_j| : c_j< 0,\; j\ \text{учтён} \}.
\end{align*}
Если множество в определении пусто, соответствующее расстояние считается бесконечным, а булев флаг наличия соседа в данном полупространстве становится ложным. Таким образом, $d_{a+}, d_{a-}, d_{c+}, d_{c-}$ --- неотрицательные расстояния до ближайшего соседа впереди/сзади вдоль цепи и влево/вправо поперёк; флаги $f_{a+}, f_{a-}, f_{c+}, f_{c-}$ фиксируют, было ли направление «занято'' наблюдением. Реализация: \code{\_nearest\_along\_cross}.

\subsection{Вспомогательные функции нелинейности и виртуального зазора}
Определим ограниченную нелинейность
\begin{equation}
\xi(x) = \tanh(x),\qquad x\in\mathbb{R}.
\end{equation}
Функция $\xi$ сжимает вклад больших расстояний и ограничивает выход интервалом $(-1,1)$, что сглаживает реакцию при тесном сближении.

Для произвольного неотрицательного $d$, константы зазора $w\ge 0$ и булева флага $f$ вводится
\begin{equation}
g(d,w,f) =
\begin{cases}
d, & f=\text{истина},\\
w, & f=\text{ложь}.
\end{cases}
\end{equation}
То есть при отсутствии соседа в заданном полупространстве вместо «бесконечного'' расстояния подставляется конечный зазор~$w$. В коде для оси вдоль при $f=\text{ложь}$ берётся $w=0$; для поперечной оси --- $w=w_{\mathrm{cross}}$, где $w_{\mathrm{cross}}>0$ (\code{W\_CROSS\_M} или \code{--w-cross}). Смысл $w_{\mathrm{cross}}$ --- имитация \textbf{виртуальной стенки} коридора с той стороны, где нет наблюдаемого соседа, чтобы скалярный сигнал не вырождался.

\subsection{Скаляры $\sigma_a$ и $\sigma_c$}
В полной динамической постановке «коридора'' в сигнал входят проекции скорости агента на оси $(\mathbf{u},\mathbf{n})$:
\begin{equation}
v_{\mathrm{along}} = v_N u_N + v_E u_E,\qquad
v_{\mathrm{cross}} = v_N n_N + v_E n_E,
\end{equation}
где $(v_N,v_E)$ --- горизонтальная скорость в NED. В рассматриваемом файле для устойчивости в SITL принято
\begin{equation}
v_{\mathrm{along}} = 0,\qquad v_{\mathrm{cross}} = 0.
\end{equation}
Тогда
\begin{align}
\label{eq:sigma}
\sigma_a &= v_{\mathrm{along}} - \xi\bigl(g(d_{a+},0,f_{a+})\bigr) + \xi\bigl(g(d_{a-},0,f_{a-})\bigr),\\
\sigma_c &= v_{\mathrm{cross}} - \xi\bigl(g(d_{c+},w_{\mathrm{cross}},f_{c+})\bigr) + \xi\bigl(g(d_{c-},w_{\mathrm{cross}},f_{c-})\bigr).
\end{align}
Величины $\sigma_a$ и $\sigma_c$ --- скалярные сигналы «распределённого'' закона в базисе цепи: они сравнивают влияние ближайших соседей с противоположных сторон по каждой оси с учётом насыщения. При $v_{\cdot}=0$ остаётся статическая нелинейная функция от положений соседей. Реализация: \code{\_sigmas\_segment\_frame}.

\subsection{Геометрические ошибки $e_{\mathrm{along}}$ и $e_{\mathrm{cross}}$}
Все агенты с известными позициями упорядочиваются по возрастанию пары $\bigl(s(p_i),\mathrm{id}_i\bigr)$, где $s$ определено в~\eqref{eq:s}. Пусть у внутреннего агента есть левый и правый сосед в этом порядке с координатами $s_L$ и $s_R$ (проекции~\eqref{eq:s} для соседей). Тогда ошибка \textbf{равномерности вдоль цепи}:
\begin{equation}
\label{eq:ealong}
e_{\mathrm{along}} = \frac{s_L + s_R}{2} - s_{\mathrm{me}},
\end{equation}
где $s_{\mathrm{me}} = s(p_{\mathrm{me}})$. Если соседа слева или справа нет (агент оказался на краю упорядоченного списка при данном снимке), полагают $e_{\mathrm{along}}=0$ в этой реализации.

Ошибка \textbf{удержания на прямой} $AB$ (signed cross-track distance):
\begin{equation}
\label{eq:ecross}
e_{\mathrm{cross}} = (N_{\mathrm{me}}-N_A)\, n_N + (E_{\mathrm{me}}-E_A)\, n_E .
\end{equation}
Положительное $e_{\mathrm{cross}}$ означает смещение в сторону $\mathbf{n}$ относительно прямой через~$A$.

Обе ошибки затем ограничиваются по модулю константами $E_{\mathrm{along}}^{\max}$ и $E_{\mathrm{cross}}^{\max}$ (\code{SPACING\_E\_ALONG\_CLAMP\_M}, \code{SPACING\_E\_CROSS\_CLAMP\_M}), чтобы ограничить максимальное «натяжение'' геометрической компоненты при выбросах измерений.

\subsection{Комбинированный вектор в плоскости North--East}
Вводятся положительные коэффициенты $k_{\mathrm{along}}$ и $k_{\mathrm{cross}}$ (\code{SPACING\_ALONG\_K}, \code{SPACING\_CROSS\_K}). Геометрическая составляющая:
\begin{equation}
\label{eq:g}
\mathbf{g}_{xy} = k_{\mathrm{along}}\, e_{\mathrm{along}}\, \mathbf{u} - k_{\mathrm{cross}}\, e_{\mathrm{cross}}\, \mathbf{n}
= (g_N,\, g_E)^\top .
\end{equation}
Знак «$-$'' перед поперечной ошибкой соответствует притягиванию к прямой в сторону, противоположную $\mathbf{n}$ при положительном $e_{\mathrm{cross}}$ (согласовано с реализацией в коде).

Вклад закона «коридора'':
\begin{equation}
\label{eq:a}
\mathbf{a}_{xy} = \sigma_a\, \mathbf{u} + \sigma_c\, \mathbf{n} = (a_N,\, a_E)^\top .
\end{equation}

Вес $w_A\in[0,1]$ (\code{ANATOLIY\_COMB\_WEIGHT}) задаёт долю «анатолиевской'' компоненты. Итоговый вектор желаемого горизонтального усилия в NED-плоскости:
\begin{equation}
\label{eq:c}
\mathbf{c}_{xy} = \mathbf{g}_{xy} + w_A\, \mathbf{a}_{xy} = (c_N,\, c_E)^\top .
\end{equation}

\subsection{Мёртвые зоны и нейтраль RC}
Пусть $\varepsilon_\sigma$ --- порог \code{SIGMA\_DEADBAND}, $\varepsilon_s$ --- порог \code{SPACING\_E\_DEADBAND\_M}. Вводится булева величина «есть геометрическое натяжение'':
\begin{equation}
\text{pull}_{\mathrm{geom}} = \bigl(|e_{\mathrm{along}}|>\varepsilon_s\bigr)\ \vee\ \bigl(|e_{\mathrm{cross}}|>\varepsilon_s\bigr).
\end{equation}
Если $\text{pull}_{\mathrm{geom}}$ ложно и одновременно $|\sigma_a|<\varepsilon_\sigma$, $|\sigma_c|<\varepsilon_\sigma$, на агент подаётся нейтральный RC (команда «не вмешиваться'').

%====================================================================
\section{Преобразование $\mathbf{c}_{xy}$ в команды RC}

\subsection{Рыскание и поворот North--East $\to$ корпус}
Пусть $\psi$ --- угол рыскания (yaw) в радианах, полученный из потока attitude; в коде к нему применяется множитель $s_y\in\{-1,+1\}$ (\code{YAW\_SIGN}) для согласования знака с конкретной сборкой SITL. Компоненты вектора $\mathbf{c}_{xy}$ в North--East обозначим $c_N, c_E$ (совпадают с $c_N,c_E$ из~\eqref{eq:c}). Желаемые «вперёд'' и «вправо'' в системе, связанной с корпусом (ось вперёд --- ось носа, вправо --- к правому крылу), до демпфирования:
\begin{equation}
\label{eq:bodyrot}
\begin{pmatrix} f_0 \\ r_0 \end{pmatrix}
=
\begin{pmatrix}
\cos\psi & \sin\psi \\
-\sin\psi & \cos\psi
\end{pmatrix}
\begin{pmatrix} c_N \\ c_E \end{pmatrix}.
\end{equation}
Здесь $f_0$ --- желаемое усилие «вперёд''; $r_0$ --- «вправо''. Реализация: \code{\_ned\_to\_body\_forward\_right}.

\subsection{Демпфирование по скорости в корпусе}
Горизонтальная скорость $(v_N,v_E)$ из стрима скоростей переводится в корпусные $(v_f,v_r)$ той же матрицей~\eqref{eq:bodyrot}. Вводятся коэффициент $k_d>0$ (\code{INTERNAL\_BODY\_V\_DAMP}) и предел насыщения $D_{\max}>0$ (\code{INTERNAL\_BODY\_V\_DAMP\_CLAMP}). Оператор $\mathrm{sat}_D$ --- ограничение по модулю: $\mathrm{sat}_D(z) = \mathrm{sign}(z)\min(|z|,D)$. Тогда
\begin{equation}
\label{eq:damp}
f = f_0 - \mathrm{sat}_{D_{\max}}(k_d v_f),\qquad
r = r_0 - \mathrm{sat}_{D_{\max}}(k_d v_r).
\end{equation}
Смысл: если аппарат уже движется в корпусных осях в ту же сторону, команду уменьшают, снижая перерегулирование.

\subsection{Пропорциональный закон, квантование PWM, минимальный шаг и slew}
Пусть $k_p>0$ --- \code{INTERNAL\_RC\_KP}, $M_{\mathrm{PWM}}$ --- \code{INTERNAL\_RC\_MAX\_PWM}, $\delta_{\min}$ --- \code{INTERNAL\_RC\_MIN\_STEP\_PWM}, $S_{\max}$ --- \code{INTERNAL\_RC\_SLEW\_PWM\_PER\_CYCLE}. Для канала pitch целевое приращение относительно нейтрали строится из~$f$; для roll используется $k_{p,r} = k_p\cdot m_r$, где $m_r$ --- \code{INTERNAL\_RC\_KP\_ROLL\_MULT} (усиление крена относительно тангажа).

На каждом шаге сначала вычисляют «сырое'' целое приращение $\mathrm{round}(k_p f)$ (и аналогично для $r$ с $k_{p,r}$), ограничивают по модулю $M_{\mathrm{PWM}}$. Если ошибка по оси ненулевая, а округление дало ноль, подставляют $\pm\delta_{\min}$ по знаку ошибки --- \textbf{минимальный шаг} (bang-bang по малым отклонениям), чтобы избежать залипания на нейтрали из-за квантования. Затем к предыдущему выходу применяют ограничение приращения за один период $T_c$ не более $S_{\max}$ (slew rate). Обозначим результаты $\Delta_{\mathrm{pitch}}$ и $\Delta_{\mathrm{roll}}$ --- целые приращения PWM после slew. Реализации: \code{\_pwm\_axis}, \code{\_slew\_int}.

\subsection{Соответствие каналам автопилота}
Пусть $U_0$ --- нейтральное значение PWM (\code{RC\_NEUTRAL}). Тогда
\begin{equation}
U_{\mathrm{pitch}} = U_0 - \Delta_{\mathrm{pitch}},\qquad
U_{\mathrm{roll}}  = U_0 + \Delta_{\mathrm{roll}},
\end{equation}
с последующим ограничением в допустимый диапазон приёмника. При несовпадении карты каналов у пользователя предусмотрен флаг \code{RC\_OVERRIDE\_SWAP\_ROLL\_PITCH}.

%====================================================================
\section{Контур якоря: слежение за точкой $A$ или $B$}

Якорь получает целевую точку $T\in\{A,B\}$ с координатами $(N_T,E_T)$. Ошибка положения в North--East:
\begin{equation}
e_N = N_T - N_{\mathrm{me}},\qquad e_E = E_T - E_{\mathrm{me}}.
\end{equation}
К компонентам применяется мёртвая зона по модулю $\varepsilon_x$ (\code{ERROR\_DEADBAND\_M}): малые $e_N,e_E$ обнуляются. Далее ошибка переводится в корпус $(e_f,e_r)$ по формуле~\eqref{eq:bodyrot} (с тем же $\psi$). Два дискретных PID-регулятора (в текущей конфигурации $K_i=K_d=0$, то есть пропорциональные каналы с разными коэффициентами) выдают $u_{\mathrm{pitch}}=\mathrm{PID}_p(e_f)$ и $u_{\mathrm{roll}}=\mathrm{PID}_r(e_r)$ с ограничением выхода \code{ENDPOINT\_OUTPUT\_LIMIT}. Маппинг на PWM совпадает по структуре с~\eqref{eq:bodyrot}--каналами для внутренних, но коэффициенты усиления другие (\code{ENDPOINT\_X\_KP} связан с продольной ошибкой в корпусе через pitch, \code{ENDPOINT\_Y\_KP} --- с поперечной через roll). После попадания в круг радиуса $\rho_{\mathrm{end}}$ (\code{ENDPOINT\_REACHED\_THRESH\_M}) якорь переключается на удержание нейтральным RC.

Реализация: \code{endpoint\_loop}, \code{\_move\_towards\_common\_with\_pid}.

%====================================================================
\section{Соответствие фрагментов кода и математических блоков}

\begin{center}
\small
\begin{tabular}{|>{\raggedright}m{4.2cm}|m{9.8cm}|}
\hline
\textbf{Функция / блок} & \textbf{Назначение} \\
\hline
\code{\_segment\_endpoints\_common\_from\_anchors} & Вычисление точек $A,B$ по~\eqref{eq:ymid}--\eqref{eq:AB}. \\
\hline
\code{\_unit\_vector\_xy}, \code{\_normal\_left\_xy} & Базис $\mathbf{u}$, $\mathbf{n}$ по~\eqref{eq:u}--\eqref{eq:n}. \\
\hline
\code{\_projection\_s} & Скаляр $s(p)$ по~\eqref{eq:s}. \\
\hline
\code{\_nearest\_along\_cross} & Расстояния $d_{a\pm},d_{c\pm}$ и флаги $f_{\cdot}$. \\
\hline
\code{\_xi}, \code{\_g}, \code{\_sigmas\_segment\_frame} & Нелинейность $\xi$, зазор $g$, сигналы $\sigma_a,\sigma_c$ по~\eqref{eq:sigma}. \\
\hline
\code{\_spacing\_errors\_projection} & Ошибки $e_{\mathrm{along}},e_{\mathrm{cross}}$ по~\eqref{eq:ealong}--\eqref{eq:ecross}. \\
\hline
\code{internal\_chain\_anatoliy\_loop} & Сборка $\mathbf{c}_{xy}$, поворот, демпфирование, PWM. \\
\hline
\code{\_pwm\_axis}, \code{\_slew\_int} & Пропорциональное квантование, минимальный шаг, slew. \\
\hline
\code{\_move\_towards\_common\_with\_pid} & Якорный PID в корпусе. \\
\hline
\code{coordinate\_exchange\_loop} & Рассылка позиций в common frame с частотой $f_{\mathrm{ex}}$. \\
\hline
\end{tabular}
\end{center}

%====================================================================
\section{Обоснование принятых решений}

\begin{description}[style=nextline,leftmargin=0cm,itemindent=0cm]
  \item[Базис $(\mathbf{u},\mathbf{n})$ привязан к отрезку $AB$.]
  Распределённый закон коридора изначально формулируется в осях «вдоль границы / поперёк''. Использование $\mathbf{u}$ и $\mathbf{n}$ делает описание \textbf{инвариантным к наклону} цепи в плоскости и согласует геометрию с законом Anatoliy без пересчёта в мировые оси симулятора.

  \item[Общая ордината $y_{\mathrm{mid}}$ для $A$ и $B$.]
  Раздельное использование $y_{\mathrm{left}}$ и $y_{\mathrm{right}}$ при реальных шумах и задержках взлёта часто даёт наклонную целевую прямую: якоря «тянут'' внутренних к разным поперечным уровням, что визуально проявляется как излом цепи. Среднее~\eqref{eq:ymid} фиксирует одну прямую для геометрической компоненты.

  \item[Сумма $\mathbf{g}_{xy} + w_A\mathbf{a}_{xy}$ с малым $w_A$.]
  Чистый «коридорный'' закон оптимизирован для движения вдоль границы с препятствиями и может конкурировать с требованием лечь на одну прямую между якорями. Доминирование $\mathbf{g}_{xy}$ обеспечивает выравнивание по проекции~$s$ и удержание на $AB$; малый вес $w_A$ добавляет нелинейную распределённую поправку без разрушения формы цепи.

  \item[Нелинейность $\tanh$.]
  Ограничивает чувствительность при малых $d_{\cdot}$ и предотвращает неограниченный рост сигнала при сближении --- стандартный приём мягкого взаимодействия в многоагентных законах.

  \item[Виртуальный зазор $w_{\mathrm{cross}}$.]
  При отсутствии наблюдаемого соседа с одной стороны поперёк подстановка конечного $w_{\mathrm{cross}}$ в $g$ сохраняет ненулевую производную по положению и имитирует наличие границы коридора.

  \item[Обнуление $v_{\mathrm{along}},v_{\mathrm{cross}}$ в~\eqref{eq:sigma}.]
  Оценки скорости из MAVLink в SITL могут быть зашумлены или несогласованы с ожидаемым земным NED; включение скорости приводило бы к дрейфу $\sigma$. Принят статический компромисс: закон опирается на положения соседей.

  \item[RC override и дискретный контур $T_c$.]
  Единый интерфейс с автопилотом в режиме удержания; минимальный шаг PWM и slew ограничивают рывки и перегрузку контура на частоте~$f_c$.

  \item[Поворот в корпус перед формированием PWM.]
  Команды roll/pitch действуют в плоскости корпуса; при $\psi\neq 0$ прямое использование $(e_N,e_E)$ или $(c_N,c_E)$ без поворота~\eqref{eq:bodyrot} исказило бы направление усилия.

  \item[\code{YAW\_SIGN} и обмен roll/pitch.]
  Устранение расхождений знака yaw между прошивкой SITL и принятой моделью NED$\to$body; учёт пользовательской карты каналов RC.
\end{description}

%====================================================================
\section{Параметры запуска и константы}

Аргументы командной строки: \code{--drones} ($N$), \code{--duration} (длительность эксперимента, 0~--- без ограничения), \code{--segment-length} ($L$), \code{--exchange-hz} ($f_{\mathrm{ex}}$), \code{--r-vis} ($R_{\mathrm{vis}}$), \code{--w-cross} ($w_{\mathrm{cross}}$). Полный перечень именованных констант и значений по умолчанию содержится в начале файла \code{linear_chain2.py}; сжатое описание на русском --- в \code{linear_chain2\_control.md}.

%====================================================================
\section*{Заключение}
В работе описан сценарий \code{linear_chain2.py}: многоагентное формирование линейной цепи в симуляторе с разделением ролей (якоря и внутренние агенты), вводом общего кадра координат, построением отрезка~$AB$ и локального базиса $(\mathbf{u},\mathbf{n})$. Для внутренних агентов получена явная форма вектора $\mathbf{c}_{xy}$ в плоскости North--East как суммы геометрической компоненты~\eqref{eq:g} и взвешенной «коридорной''~\eqref{eq:a}, далее показан переход к командам RC через поворот в корпус~\eqref{eq:bodyrot}, демпфирование~\eqref{eq:damp} и дискретное ограниченное преобразование в PWM. Приведены обоснования ключевых компромиссов (средняя ордината концов, гибридный закон, $v=0$ в $\sigma$, квантование и slew). Для проверки соответствия формул кодовой базе рекомендуется сопоставлять номера уравнений с указанными функциями при каждом изменении исходного файла.

\end{document}
